% Options for packages loaded elsewhere
\PassOptionsToPackage{unicode}{hyperref}
\PassOptionsToPackage{hyphens}{url}
%
\documentclass[
]{article}
\usepackage{amsmath,amssymb}
\usepackage{iftex}
\ifPDFTeX
  \usepackage[T1]{fontenc}
  \usepackage[utf8]{inputenc}
  \usepackage{textcomp} % provide euro and other symbols
\else % if luatex or xetex
  \usepackage{unicode-math} % this also loads fontspec
  \defaultfontfeatures{Scale=MatchLowercase}
  \defaultfontfeatures[\rmfamily]{Ligatures=TeX,Scale=1}
\fi
\usepackage{lmodern}
\ifPDFTeX\else
  % xetex/luatex font selection
\fi
% Use upquote if available, for straight quotes in verbatim environments
\IfFileExists{upquote.sty}{\usepackage{upquote}}{}
\IfFileExists{microtype.sty}{% use microtype if available
  \usepackage[]{microtype}
  \UseMicrotypeSet[protrusion]{basicmath} % disable protrusion for tt fonts
}{}
\makeatletter
\@ifundefined{KOMAClassName}{% if non-KOMA class
  \IfFileExists{parskip.sty}{%
    \usepackage{parskip}
  }{% else
    \setlength{\parindent}{0pt}
    \setlength{\parskip}{6pt plus 2pt minus 1pt}}
}{% if KOMA class
  \KOMAoptions{parskip=half}}
\makeatother
\usepackage{xcolor}
\usepackage{longtable,booktabs,array}
\usepackage{calc} % for calculating minipage widths
% Correct order of tables after \paragraph or \subparagraph
\usepackage{etoolbox}
\makeatletter
\patchcmd\longtable{\par}{\if@noskipsec\mbox{}\fi\par}{}{}
\makeatother
% Allow footnotes in longtable head/foot
\IfFileExists{footnotehyper.sty}{\usepackage{footnotehyper}}{\usepackage{footnote}}
\makesavenoteenv{longtable}
\usepackage{graphicx}
\makeatletter
\def\maxwidth{\ifdim\Gin@nat@width>\linewidth\linewidth\else\Gin@nat@width\fi}
\def\maxheight{\ifdim\Gin@nat@height>\textheight\textheight\else\Gin@nat@height\fi}
\makeatother
% Scale images if necessary, so that they will not overflow the page
% margins by default, and it is still possible to overwrite the defaults
% using explicit options in \includegraphics[width, height, ...]{}
\setkeys{Gin}{width=\maxwidth,height=\maxheight,keepaspectratio}
% Set default figure placement to htbp
\makeatletter
\def\fps@figure{htbp}
\makeatother
\setlength{\emergencystretch}{3em} % prevent overfull lines
\providecommand{\tightlist}{%
  \setlength{\itemsep}{0pt}\setlength{\parskip}{0pt}}
\setcounter{secnumdepth}{-\maxdimen} % remove section numbering
\ifLuaTeX
  \usepackage{selnolig}  % disable illegal ligatures
\fi
\IfFileExists{bookmark.sty}{\usepackage{bookmark}}{\usepackage{hyperref}}
\IfFileExists{xurl.sty}{\usepackage{xurl}}{} % add URL line breaks if available
\urlstyle{same}
\hypersetup{
  hidelinks,
  pdfcreator={LaTeX via pandoc}}

\author{}
\date{}

\begin{document}

\textbf{MNIST Digit Classification Using Federated CNN with
Nature-Inspired Feature Selection}

\emph{A}

\emph{\textbf{Project Report}}

\emph{submitted in partial fulfillment of the}

\emph{requirements for the award of the degree of}

\textbf{BACHELOR OF TECHNOLOGY}

in

\textbf{COMPUTER SCIENCE \& ENGINEERING}

by

\begin{longtable}[]{@{}
  >{\raggedright\arraybackslash}p{(\columnwidth - 2\tabcolsep) * \real{0.5333}}
  >{\raggedright\arraybackslash}p{(\columnwidth - 2\tabcolsep) * \real{0.4667}}@{}}
\toprule\noalign{}
\begin{minipage}[b]{\linewidth}\raggedright
\textbf{Name}
\end{minipage} & \begin{minipage}[b]{\linewidth}\raggedright
\textbf{Roll No.}
\end{minipage} \\
\midrule\noalign{}
\endhead
\bottomrule\noalign{}
\endlastfoot
Varchasva Choudhary & R2142231573 \\
Naman Dixit & R2142231696 \\
Charvi Agarwal & R2142231186 \\
Kartik Goyal & R2142231215 \\
\end{longtable}

\emph{under the guidance of}

\textbf{Narita Sarkar}

\includegraphics[width=1.29167in,height=1.29167in]{media/image1.jpeg}

\textbf{School of Computer Science}

\textbf{University of Petroleum \& Energy Studies}

\textbf{Bidholi, Via Prem Nagar, Dehradun, Uttarakhand}

\textbf{May -- 2025}

\hypertarget{candidates-declaration}{%
\section{CANDIDATE\textquotesingle S
DECLARATION}\label{candidates-declaration}}

We hereby certify that the project work entitled \textbf{"MNIST Digit
Classification Using Federated CNN with Nature-Inspired Feature
Selection"} in partial fulfilment of the requirements for the award of
the Degree of BACHELOR OF TECHNOLOGY in COMPUTER SCIENCE AND ENGINEERING
and submitted to the Department of Systemics, School of Computer
Science, University of Petroleum \& Energy Studies, Dehradun, is an
authentic record of our work carried out during the period from January,
2025 to May, 2025 under the supervision of \textbf{Narita Sarkar mam.}

The matter presented in this project has not been submitted by us for
the award of any other degree of this or any other University.

\textbf{(Varchasva Choudhary, Naman Dixit, Charvi Agarwal, Kartik
Goyal)}

This is to certify that the above statement made by the candidates is
correct to the best of my knowledge.

Date: 5 May 2025 \textbf{Narita Sarkar}

\emph{Project Guide}

\hypertarget{acknowledgement}{%
\section{ACKNOWLEDGEMENT}\label{acknowledgement}}

We wish to express our deep gratitude to our guide {[}Name{]}, for all
the advice, encouragement, and constant support provided throughout this
project work. This work would not have been possible without their
invaluable guidance and constructive suggestions.

We sincerely thank {[}Name of HoD{]}, Head of the Department of Computer
Science, for the great support extended in carrying out this project and
for providing all the necessary infrastructure and resources.

We are also grateful to the Dean, School of Computer Science, UPES, for
providing the necessary facilities to carry out our project work
successfully. We also thank our Course Coordinator and Activity
Coordinator for providing timely support and information during the
completion of this project.

We would like to thank all our friends for their help and constructive
criticism during our project work. Finally, we have no words to express
our sincere gratitude to our parents for their unconditional support and
encouragement throughout.

\begin{longtable}[]{@{}
  >{\raggedright\arraybackslash}p{(\columnwidth - 8\tabcolsep) * \real{0.1994}}
  >{\raggedright\arraybackslash}p{(\columnwidth - 8\tabcolsep) * \real{0.1994}}
  >{\raggedright\arraybackslash}p{(\columnwidth - 8\tabcolsep) * \real{0.1994}}
  >{\raggedright\arraybackslash}p{(\columnwidth - 8\tabcolsep) * \real{0.1994}}
  >{\raggedright\arraybackslash}p{(\columnwidth - 8\tabcolsep) * \real{0.2023}}@{}}
\toprule\noalign{}
\begin{minipage}[b]{\linewidth}\raggedright
\textbf{Name}
\end{minipage} & \begin{minipage}[b]{\linewidth}\raggedright
Varchasva Choudhary
\end{minipage} & \begin{minipage}[b]{\linewidth}\raggedright
Naman Dixit
\end{minipage} & \begin{minipage}[b]{\linewidth}\raggedright
Charvi Agarwal
\end{minipage} & \begin{minipage}[b]{\linewidth}\raggedright
Kartik Goyal
\end{minipage} \\
\midrule\noalign{}
\endhead
\bottomrule\noalign{}
\endlastfoot
\textbf{Roll No.} & R2142231573 & R2142231696 & R2142231186 &
R2142231215 \\
\end{longtable}

\hypertarget{abstract}{%
\section{ABSTRACT}\label{abstract}}

This project presents a comprehensive study of MNIST handwritten digit
classification using a Federated Convolutional Neural Network (FedCNN)
integrated with nature-inspired wrapper-based feature selection
algorithms. The system combines privacy-preserving federated learning
with intelligent feature selection to achieve high classification
accuracy with reduced feature dimensionality.

The Federated CNN architecture distributes training across five
simulated clients using a Reinforcement Learning (RL)-based Q-Agent for
intelligent client selection in each communication round. After
federated training, the global model extracts 128-dimensional deep
feature representations from the MNIST dataset. These features are then
subjected to three nature-inspired metaheuristic optimization
algorithms: Genetic Algorithm (GenA\_Conv), Grey Wolf Optimizer
(GreyW\_Conv), and Particle Swarm Optimization (ParticleS\_Conv), which
serve as wrapper-based feature selectors to identify the most
discriminative subset of CNN features.

The system is evaluated across five different train-test splits (65:35,
70:30, 75:25, 80:20, 85:15) and benchmarked against a Baseline Logistic
Regression model (using raw 784-pixel features) and Conv\_NN (using all
128 CNN features without selection). Performance metrics including
Accuracy, Precision, Recall, and F1-Score are compared using heatmaps,
convergence plots, confusion matrices, correlation heatmaps, and
SHAP-based interpretability analysis.

Results demonstrate that nature-inspired feature selection methods,
particularly PSO and GWO, achieve competitive or superior classification
accuracy compared to the full feature set while significantly reducing
feature dimensionality --- improving model interpretability, efficiency,
and generalization.

\hypertarget{table-of-contents}{%
\section{TABLE OF CONTENTS}\label{table-of-contents}}

\textbf{S.No. Contents Page No.}

1. Introduction 1

1.1 Background and Motivation 1

1.2 Problem Statement 2

1.3 Objectives 2

2. System Analysis 3

2.1 Existing Systems 3

2.2 Proposed System 4

2.3 Modules Description 5

3. Design 6

3.1 System Architecture 6

3.2 FedCNN Architecture 7

3.3 Optimizer Designs 8

4. Implementation 9

4.1 Data Preprocessing 9

4.2 Federated Training 10

4.3 Feature Extraction \& Selection 11

5. Results and Discussion 12

5.1 Performance Metrics 12

5.2 Comparative Analysis 13

5.3 Best Method Selection 14

6. Conclusion and Future Work 15

References 16

\hypertarget{introduction}{%
\section{1. INTRODUCTION}\label{introduction}}

\hypertarget{background-and-motivation}{%
\subsection{1.1 Background and
Motivation}\label{background-and-motivation}}

Handwritten digit recognition remains a fundamental benchmark problem in
the field of machine learning and computer vision. The MNIST dataset,
comprising 70,000 grayscale images of handwritten digits (0--9) at 28×28
pixel resolution, has served as the standard testbed for evaluating
classification algorithms for decades. While conventional deep learning
approaches achieve near-human accuracy on MNIST, modern challenges
demand solutions that simultaneously address privacy concerns, model
interpretability, and computational efficiency.

Federated Learning (FL) has emerged as a privacy-preserving distributed
machine learning paradigm wherein model training is performed across
multiple decentralized clients without sharing raw data. This is
particularly relevant in healthcare, finance, and mobile applications,
where data privacy is paramount. Combining FL with deep Convolutional
Neural Networks (CNNs) leverages the feature extraction power of CNNs
while preserving the privacy guarantees of federated training.

Feature selection is a critical step in machine learning pipelines.
Nature-inspired metaheuristic algorithms --- including Genetic
Algorithms (GA), Grey Wolf Optimizer (GWO), and Particle Swarm
Optimization (PSO) --- have demonstrated remarkable effectiveness as
wrapper-based feature selectors. These methods treat feature selection
as a combinatorial optimization problem, searching for the optimal
feature subset that maximizes classification performance while
minimizing the number of selected features.

\hypertarget{problem-statement}{%
\subsection{1.2 Problem Statement}\label{problem-statement}}

Standard CNN-based approaches for MNIST classification train models on
centralized data, which raises privacy concerns in real-world
deployments. Moreover, raw CNN feature representations (e.g.,
128-dimensional embeddings) may contain redundant or irrelevant
dimensions that add computational overhead without contributing to
classification accuracy. There exists a need for a unified framework
that:

\begin{itemize}
\item
  Enables privacy-preserving distributed training across multiple
  clients.
\item
  Intelligently selects the most informative CNN feature dimensions.
\item
  Benchmarks multiple optimization strategies across varying train-test
  splits.
\item
  Provides interpretable insights into which features drive
  classification decisions.
\end{itemize}

\hypertarget{objectives}{%
\subsection{1.3 Objectives}\label{objectives}}

The primary objectives of this project are:

\begin{itemize}
\item
  To design and implement a Federated CNN (FedCNN) for MNIST digit
  classification using a Q-Agent for RL-based client selection.
\item
  To extract 128-dimensional deep feature embeddings from the trained
  global FedCNN model.
\item
  To implement three nature-inspired wrapper feature selectors: Genetic
  Algorithm (GenA\_Conv), Grey Wolf Optimizer (GreyW\_Conv), and
  Particle Swarm Optimization (ParticleS\_Conv).
\item
  To evaluate all four methods --- Conv\_NN, GenA\_Conv, GreyW\_Conv,
  ParticleS\_Conv --- across five train-test splits using Accuracy,
  Precision, Recall, and F1-Score.
\item
  To perform SHAP-based interpretability analysis on the best-performing
  method.
\item
  To identify the optimal feature selection strategy through composite
  scoring.
\end{itemize}

\hypertarget{system-analysis}{%
\section{2. SYSTEM ANALYSIS}\label{system-analysis}}

\hypertarget{existing-systems}{%
\subsection{2.1 Existing Systems}\label{existing-systems}}

Traditional approaches to MNIST classification include Support Vector
Machines (SVM), k-Nearest Neighbors (k-NN), and simple Multilayer
Perceptrons (MLP). While these methods achieved reasonable accuracy,
they require centralized data storage and use raw pixel features, which
are high-dimensional and redundant.

Deep learning approaches --- particularly CNNs --- have pushed MNIST
accuracy beyond 99\%. However, standard CNN training centralizes all
client data, making it unsuitable for privacy-sensitive applications.
Feature selection methods such as Principal Component Analysis (PCA)
reduce dimensionality but are unsupervised and may discard
class-discriminative information.

Existing federated learning frameworks (e.g., FedAvg) aggregate model
weights from clients but do not incorporate intelligent client selection
or combine feature optimization with federated training. The limitations
of existing systems motivated the development of the proposed integrated
framework.

\hypertarget{proposed-system}{%
\subsection{2.2 Proposed System}\label{proposed-system}}

The proposed system integrates three major components: Federated CNN
training with RL-based client selection, CNN feature extraction, and
nature-inspired wrapper feature selection. The workflow proceeds as
follows:

First, the MNIST dataset is preprocessed using normalization,
standardization, and threshold-based denoising. The data is divided into
five simulated federated clients. A Q-Agent selects the top-3 clients
per federated round based on validation accuracy rewards. Local CNN
models are trained for one epoch per client, and their weights are
aggregated using Federated Averaging (FedAvg). This process repeats for
five federated rounds.

After federated training, the global FedCNN model serves as a feature
extractor, producing 128-dimensional embeddings per sample. These
embeddings are then input to three nature-inspired optimizers (GA, GWO,
PSO), each acting as a wrapper to select the most informative feature
subset. A Logistic Regression classifier is trained on the selected
features for each method. Finally, all methods are evaluated across five
train-test splits using standard classification metrics.

\hypertarget{modules-description}{%
\subsection{2.3 Modules Description}\label{modules-description}}

\hypertarget{data-preprocessing-module}{%
\subsubsection{2.3.1 Data Preprocessing
Module}\label{data-preprocessing-module}}

The MNIST dataset (70,000 samples, 784 features each) is fetched using
scikit-learn\textquotesingle s fetch\_openml API. Pixel values are
normalized to {[}0, 1{]}, standardized using z-score normalization, and
binarized using a threshold of 0.3 to reduce noise. This preprocessing
pipeline ensures consistent input distributions across all experiments.

\hypertarget{federated-cnn-training-module}{%
\subsubsection{2.3.2 Federated CNN Training
Module}\label{federated-cnn-training-module}}

The FedCNN architecture consists of two convolutional blocks (Conv2D →
BatchNorm → ReLU → MaxPool), followed by a 128-unit dense feature layer
with Dropout(0.5) and a 10-class output layer. Training is distributed
across 5 clients with a Q-Agent selecting 3 clients per round based on
validation accuracy rewards. FedAvg aggregates client weights after each
round.

\hypertarget{feature-extraction-module}{%
\subsubsection{2.3.3 Feature Extraction
Module}\label{feature-extraction-module}}

The trained global FedCNN model is used to extract 128-dimensional
feature vectors from train and test sets. These CNN features serve as
the input space for subsequent classifier training and
optimization-based feature selection.

\hypertarget{nature-inspired-feature-selection-module}{%
\subsubsection{2.3.4 Nature-Inspired Feature Selection
Module}\label{nature-inspired-feature-selection-module}}

Three metaheuristic algorithms are implemented as wrapper feature
selectors. Each algorithm evaluates candidate feature subsets by
training a Logistic Regression classifier and measuring validation
accuracy as the fitness score. The Genetic Algorithm (GA) uses
selection, crossover, and mutation operators. Grey Wolf Optimizer (GWO)
mimics the hierarchical hunting behaviour of grey wolves. Particle Swarm
Optimization (PSO) models social behaviour of particles converging on
global optima.

\hypertarget{design}{%
\section{3. DESIGN}\label{design}}

\hypertarget{system-architecture}{%
\subsection{3.1 System Architecture}\label{system-architecture}}

The overall system architecture follows a sequential pipeline: raw MNIST
data → preprocessing → federated client distribution → RL-based client
selection → local CNN training → FedAvg aggregation → global model
evaluation → 128-D feature extraction → nature-inspired feature
selection (GA / GWO / PSO) → Logistic Regression classification →
multi-split evaluation → SHAP interpretability.

The federated component ensures data privacy by never sharing raw
training data between clients or with the central server. Only model
weight updates are communicated. The RL-based Q-Agent adds an adaptive
layer that preferentially selects high-performing clients in subsequent
rounds, improving convergence efficiency.

\hypertarget{fedcnn-architecture}{%
\subsection{3.2 FedCNN Architecture}\label{fedcnn-architecture}}

The FedCNN consists of the following layers:

\begin{longtable}[]{@{}
  >{\raggedright\arraybackslash}p{(\columnwidth - 4\tabcolsep) * \real{0.3324}}
  >{\raggedright\arraybackslash}p{(\columnwidth - 4\tabcolsep) * \real{0.3324}}
  >{\raggedright\arraybackslash}p{(\columnwidth - 4\tabcolsep) * \real{0.3353}}@{}}
\toprule\noalign{}
\begin{minipage}[b]{\linewidth}\raggedright
\textbf{Layer}
\end{minipage} & \begin{minipage}[b]{\linewidth}\raggedright
\textbf{Configuration}
\end{minipage} & \begin{minipage}[b]{\linewidth}\raggedright
\textbf{Output Shape}
\end{minipage} \\
\midrule\noalign{}
\endhead
\bottomrule\noalign{}
\endlastfoot
Conv2D + BN + ReLU & 32 filters, 3×3 kernel, padding=1 & 32 × 28 × 28 \\
MaxPool2D & 2×2 & 32 × 14 × 14 \\
Conv2D + BN + ReLU & 64 filters, 3×3 kernel, padding=1 & 64 × 14 × 14 \\
MaxPool2D & 2×2 & 64 × 7 × 7 \\
Flatten & --- & 3136 \\
Linear + ReLU + Dropout & 128 units, p=0.5 & 128 (feature vector) \\
Linear (classifier) & 10 classes & 10 \\
\end{longtable}

\hypertarget{optimizer-designs}{%
\subsection{3.3 Optimizer Designs}\label{optimizer-designs}}

\hypertarget{genetic-algorithm-gena_conv}{%
\subsubsection{3.3.1 Genetic Algorithm
(GenA\_Conv)}\label{genetic-algorithm-gena_conv}}

The GA encodes candidate feature subsets as binary chromosomes of length
128. Fitness is measured as Logistic Regression accuracy on the
validation set. The GA applies tournament selection, single-point
crossover, and bit-flip mutation over multiple generations, maintaining
a population of candidate solutions that evolve towards optimal feature
subsets.

\hypertarget{grey-wolf-optimizer-greyw_conv}{%
\subsubsection{3.3.2 Grey Wolf Optimizer
(GreyW\_Conv)}\label{grey-wolf-optimizer-greyw_conv}}

GWO models the social hierarchy of grey wolves: alpha (best solution),
beta (second-best), and delta (third-best) guide the search. Candidate
wolves (solutions) update their positions based on the influence of
alpha, beta, and delta, enabling balanced exploration and exploitation
in the feature selection search space.

\hypertarget{particle-swarm-optimization-particles_conv}{%
\subsubsection{3.3.3 Particle Swarm Optimization
(ParticleS\_Conv)}\label{particle-swarm-optimization-particles_conv}}

PSO maintains a swarm of particles, each representing a candidate
feature subset. Each particle updates its velocity and position based on
its personal best and the global best solution found by the swarm. The
inertia weight and cognitive/social coefficients are tuned to balance
exploration and exploitation, converging to the globally optimal feature
subset.

\hypertarget{implementation}{%
\section{4. IMPLEMENTATION}\label{implementation}}

\hypertarget{data-preprocessing}{%
\subsection{4.1 Data Preprocessing}\label{data-preprocessing}}

The MNIST dataset was fetched using scikit-learn\textquotesingle s
fetch\_openml API, yielding X ∈ ℝ\^{}(70000×784) and labels y ∈
\{0,...,9\}. The preprocessing pipeline consists of three steps:

\begin{itemize}
\item
  Normalization: X\_scaled = X / 255.0, mapping pixel intensities to
  {[}0, 1{]}.
\item
  Standardization: X\_scaled = (X\_scaled − μ) / σ, where μ and σ are
  the global mean and standard deviation.
\item
  Denoising: Binary thresholding at 0.3 --- pixels with standardized
  value \textgreater{} 0.3 are set to 1.0, others to 0.0.
\end{itemize}

Five train-test splits were defined: 65:35, 70:30, 75:25, 80:20, and
85:15. Stratified splitting ensures class balance across splits. All
experiments use random\_state=42 for reproducibility.

\hypertarget{federated-training}{%
\subsection{4.2 Federated Training}\label{federated-training}}

The federated training pipeline proceeds as follows:

\begin{itemize}
\item
  The training set (80:20 split) is partitioned into 5 equal-sized
  client datasets using non-overlapping contiguous segments.
\item
  A Q-Agent (ClientSelectorQAgent) maintains a Q-table of size 5
  initialized to zeros. In each federated round, the top-3 clients are
  selected based on descending Q-table values.
\item
  Selected clients train local deep-copies of the global FedCNN model
  for 1 epoch using Adam optimizer (lr=0.001) and CrossEntropyLoss.
\item
  Local model performance on a held-out validation set (15\% of training
  data) provides the reward signal for Q-table updates.
\item
  FedAvg computes the element-wise mean of client state dictionaries to
  update the global model.
\item
  This process repeats for 5 federated rounds, after which the global
  model achieves stable classification performance.
\end{itemize}

\hypertarget{feature-extraction-and-selection}{%
\subsection{4.3 Feature Extraction and
Selection}\label{feature-extraction-and-selection}}

The trained global FedCNN model is used to extract 128-dimensional
feature embeddings from all train and test samples using the
feature\_layer output (before the final classifier). These CNN features
capture high-level visual representations learned during federated
training.

The three nature-inspired optimizers are then applied as wrapper
selectors on the 128-dimensional feature space:

\begin{itemize}
\item
  GenA\_Conv (GA): Binary chromosome representation; fitness = LR
  validation accuracy; operators: tournament selection, single-point
  crossover, bit-flip mutation.
\item
  GreyW\_Conv (GWO): Alpha-beta-delta hierarchy; position updates model
  grey wolf pack hunting convergence on optimal feature subsets.
\item
  ParticleS\_Conv (PSO): Swarm of particles; velocity and position
  updates based on personal best and global best; inertia weight
  linearly decreasing from 0.9 to 0.4.
\end{itemize}

For each optimizer, the selected feature indices are used to train a
final Logistic Regression classifier (max\_iter=1000, lbfgs solver),
which is evaluated on the held-out test set. Experiments are repeated
across all five train-test splits.

\hypertarget{results-and-discussion}{%
\section{5. RESULTS AND DISCUSSION}\label{results-and-discussion}}

\hypertarget{performance-metrics}{%
\subsection{5.1 Performance Metrics}\label{performance-metrics}}

All methods are evaluated using four standard classification metrics:
Accuracy, Precision (weighted), Recall (weighted), and F1-Score
(weighted). These metrics are computed using
scikit-learn\textquotesingle s accuracy\_score, precision\_score,
recall\_score, and f1\_score functions with zero\_division=0 and
average=\textquotesingle weighted\textquotesingle{} to handle class
imbalance appropriately.

\hypertarget{comparative-analysis}{%
\subsection{5.2 Comparative Analysis}\label{comparative-analysis}}

The following observations are drawn from the multi-split comparative
evaluation:

\hypertarget{baseline-logistic-regression}{%
\subsubsection{5.2.1 Baseline Logistic
Regression}\label{baseline-logistic-regression}}

The baseline Logistic Regression, trained on all 784 raw pixel features,
serves as the lower-bound reference. Despite using all available input
information, the linear model struggles to capture non-linear digit
patterns inherent in raw pixel space, yielding modest accuracy across
all splits.

\hypertarget{conv_nn-all-128-cnn-features}{%
\subsubsection{5.2.2 Conv\_NN (All 128 CNN
Features)}\label{conv_nn-all-128-cnn-features}}

Training Logistic Regression on all 128 CNN features extracted from the
FedCNN substantially improves over the baseline, confirming that
federated deep feature extraction captures discriminative visual
representations. Accuracy remains consistent across train-test splits,
indicating low variance.

\hypertarget{gena_conv-greyw_conv-particles_conv}{%
\subsubsection{5.2.3 GenA\_Conv, GreyW\_Conv,
ParticleS\_Conv}\label{gena_conv-greyw_conv-particles_conv}}

All three nature-inspired selectors achieve competitive accuracy to
Conv\_NN while selecting a fraction of the 128 available features. GWO
and PSO tend to achieve the highest fitness scores with fewer features,
demonstrating superior exploration-exploitation balance. The GA
converges reliably but may select slightly more features to achieve
similar accuracy. Convergence plots illustrate that all three optimizers
show monotonically increasing fitness trends, confirming proper
convergence behaviour.

Confusion matrices across all four methods confirm that digit classes
such as \textquotesingle1\textquotesingle{} and
\textquotesingle0\textquotesingle{} are classified with highest
precision, while \textquotesingle4\textquotesingle{} and
\textquotesingle9\textquotesingle{} exhibit the most inter-class
confusion --- consistent with their visual similarity.

\hypertarget{best-method-selection}{%
\subsection{5.3 Best Method Selection}\label{best-method-selection}}

A composite scoring methodology is used to objectively rank all four
methods. The composite score is defined as:

\emph{\textbf{Composite Score = 0.40 × Mean Accuracy + 0.40 × Mean
F1-Score + 0.20 × Feature Reduction Ratio}}

where Feature Reduction Ratio = 1 − (mean selected features / 128). This
composite score rewards both classification performance and feature
economy. The method achieving the highest composite score across all
five splits is identified as the overall winner.

SHAP (SHapley Additive exPlanations) analysis is performed on the best
method to identify the top-20 most influential CNN features driving
classification decisions. The SHAP beeswarm and bar plots provide
interpretable feature importance rankings, enabling transparency in
model decision-making.

\hypertarget{conclusion-and-future-work}{%
\section{6. CONCLUSION AND FUTURE
WORK}\label{conclusion-and-future-work}}

This project successfully developed and evaluated an integrated
framework combining Federated CNN training with RL-based client
selection and nature-inspired wrapper feature selection for MNIST digit
classification. The system demonstrates that:

\begin{itemize}
\item
  Federated Learning effectively trains a high-accuracy CNN model
  without centralizing client data, preserving data privacy.
\item
  RL-based Q-Agent client selection improves federated convergence by
  dynamically prioritizing high-performing clients.
\item
  Nature-inspired feature selectors (GA, GWO, PSO) successfully identify
  compact, discriminative subsets of CNN features, reducing
  dimensionality while maintaining or improving classification accuracy.
\item
  The composite scoring methodology provides an objective framework for
  selecting the optimal feature selection strategy in multi-criteria
  evaluation scenarios.
\item
  SHAP analysis enhances model interpretability by identifying the most
  influential features driving digit classification decisions.
\end{itemize}

The proposed framework achieves a strong balance between
privacy-preservation, classification accuracy, model interpretability,
and computational efficiency --- addressing key limitations of existing
centralized deep learning approaches.

\hypertarget{future-work}{%
\subsection{Future Work}\label{future-work}}

Potential extensions of this work include:

\begin{itemize}
\item
  Applying the framework to real-world federated scenarios with
  heterogeneous (non-IID) client data distributions.
\item
  Exploring deep feature selection using neural architecture search
  (NAS) instead of linear wrappers.
\item
  Extending to multi-modal datasets beyond MNIST (e.g., CIFAR-10,
  medical imaging datasets).
\item
  Incorporating differential privacy mechanisms to provide formal
  privacy guarantees during federated training.
\item
  Investigating online feature selection where the optimal feature
  subset evolves dynamically across federated rounds.
\end{itemize}

\hypertarget{references}{%
\section{REFERENCES}\label{references}}

(As per their appearance in the chapters)

{[}1{]} LeCun, Y., Bottou, L., Bengio, Y., \& Haffner, P. (1998).
Gradient-based learning applied to document recognition. Proceedings of
the IEEE, 86(11), 2278--2324.

{[}2{]} McMahan, H. B., Moore, E., Ramage, D., Hampson, S., \& y Arcas,
B. A. (2017). Communication-efficient learning of deep networks from
decentralized data. AISTATS 2017.

{[}3{]} Holland, J. H. (1975). Adaptation in Natural and Artificial
Systems. University of Michigan Press.

{[}4{]} Mirjalili, S., Mirjalili, S. M., \& Lewis, A. (2014). Grey Wolf
Optimizer. Advances in Engineering Software, 69, 46--61.

{[}5{]} Kennedy, J., \& Eberhart, R. (1995). Particle Swarm
Optimization. Proceedings of ICNN\textquotesingle95, 4, 1942--1948.

{[}6{]} Lundberg, S. M., \& Lee, S. I. (2017). A unified approach to
interpreting model predictions. NeurIPS 30.

{[}7{]} Paszke, A., et al. (2019). PyTorch: An imperative style,
high-performance deep learning library. NeurIPS 32.

{[}8{]} Pedregosa, F., et al. (2011). Scikit-learn: Machine learning in
Python. Journal of Machine Learning Research, 12, 2825--2830.

\end{document}
